% MMW-METHOD-CONTRACT: ead1dc31de5453e83924fd80c141e14521697c804489a275435dbb5338b47dac
% MMW-ID: OBJ-Q1
% MMW-ID: OBJ-Q2
% MMW-ID: OBJ-Q3
% MMW-ID: OBJ-Q4
% MMW-ID: CON-Q1-1
% MMW-ID: CON-Q1-2
% MMW-ID: CON-Q1-3
% MMW-ID: CON-Q1-4
% MMW-ID: CON-Q1-5
% MMW-ID: CON-Q1-6
% MMW-ID: CON-Q1-7
% MMW-ID: CON-Q2-1
% MMW-ID: CON-Q2-2
% MMW-ID: CON-Q2-3
% MMW-ID: CON-Q2-4
% MMW-ID: CON-Q2-5
% MMW-ID: CON-Q2-6
% MMW-ID: CON-Q3-1
% MMW-ID: CON-Q3-2
% MMW-ID: CON-Q3-3
% MMW-ID: CON-Q3-4
% MMW-ID: CON-Q3-5
% MMW-ID: CON-Q3-6
% MMW-ID: CON-Q3-7
% MMW-ID: CON-Q4-1
% MMW-ID: CON-Q4-2
% MMW-ID: CON-Q4-3
% MMW-ID: CON-Q4-4
% MMW-ID: CON-Q4-5
% MMW-ID: CON-Q4-6
% MMW-ALGORITHM: 严格闭环预测；五方案增量任务局部搜索；离散SOC动态规划；按情景动态重建邻域的任务-储能交替有限搜索
\section{模型的建立与求解}

\subsection{统一时间、资源与能源口径}
任务 $i$ 的 GPU 数、到达时刻、持续小时数和有效截止分别记为 $g_i,a_i,h_i,F_i$，其中 $h_i=d_i/60$，$F_i=\min\{l_i,2406\}$。任务在整数时刻 $s_i$ 开工，与小时区间 $[t,t+1)$ 的重叠量为
\begin{equation}
\omega_{it}=\max\{0,\min(t+1,s_i+h_i)-\max(t,s_i)\},\qquad
\chi_{it}=\mathbf 1\{\omega_{it}>0\}.
\end{equation}
$\chi_{it}$ 用于瞬时 GPU 与功率边界，$\omega_{it}$ 用于电量。合法候选集合满足
\begin{equation}
\Omega_i=\{(r,s):s\ge a_i,\ s+h_i\le F_i,\ \ell_{o_i r}\le \ell_i^{\max},\ \omega_{i,2406}=0\}.
\end{equation}
每个任务恰选一个候选，实时任务还满足 $s_i=a_i$。区域 $r$ 在小时 $t$ 的 GPU、IT 与设施功率约束为
\begin{align}
\sum_i g_i\chi_{it}\mathbf 1(r_i=r)&\le \bar G_r,\\
L_{rt}^{N}+\sum_i g_i p_{k_i}^{GPU}\chi_{it}\mathbf 1(r_i=r)&\le \bar P_r^{IT},\\
\pi_r\!\left(L_{rt}^{N}+\sum_i g_i p_{k_i}^{GPU}\chi_{it}\mathbf 1(r_i=r)\right)&\le \bar P_r^{F}.
\end{align}
AI 平均负荷则以 $\omega_{it}$ 累计。运行成本、购电碳排放与新能源利用率统一为
\begin{align}
C&=\sum_{r,t}(c^{buy}_{rt}G^{buy}_{rt}-c^{sell}_{rt}G^{sell}_{rt}),\\
E&=\sum_{r,t}\kappa_{rt}G^{buy}_{rt},\\
U&=\frac{\sum_{r,t}(R^{use}_{rt}+R^{ch}_{rt}+R^{sell}_{rt})}{\sum_{r,t}R^{av}_{rt}},
\end{align}
其中 $U$ 仅在分母为正时计算。网络时延按任务 GPU-hour 加权，服务质量由实时 SLA、按期完成率和等待统计共同解释。

\subsection{问题一：严格闭环预测与基础调度}
\subsubsection{候选预测模型与选模}
对区域 $r$、任务类型 $k$ 的到达 GPU 序列 $D_{rkt}$，Huber 候选使用滞后 $1,2,3,24,48,168$ 及日、周正余弦特征。其损失为
\begin{equation}
\min_{b,\theta,\sigma>0}n\sigma+\sigma\sum_t\rho_{\epsilon}\!\left(\frac{D_{rkt}-b-x_t^\mathsf T\theta}{\sigma}\right)+\alpha\|\theta\|_2^2.
\end{equation}
其特征向量明确写为
\begin{equation}
\begin{aligned}
x_{rk,t}=\big[&D_{rk,t-1},D_{rk,t-2},D_{rk,t-3},D_{rk,t-24},D_{rk,t-48},D_{rk,t-168},\\
&\sin(2\pi t/24),\cos(2\pi t/24),\sin(2\pi t/168),\cos(2\pi t/168)\big]^{\mathsf T}.
\end{aligned}
\end{equation}
短滞后捕捉临近惯性，24 小时与 168 小时滞后分别对应日、周重复结构，正余弦项避免把周期边界处理成跳变。

超参数网格为 $\epsilon\in\{1.10,1.35,1.75\}$、$\alpha\in\{10^{-4},10^{-3},10^{-2}\}$，并与 24 小时、168 小时季节朴素模型比较。第 0 至 2351 小时训练，第 2352 至 2375 小时验证，第 2376 至 2399 小时测试。验证阶段按 RMSE、MAE、WAPE 和固定模型次序破除并列；测试阶段闭环递推，字段 \texttt{FeatureUsesTestTruth} 全为 false。

\begin{table}[H]\centering
\caption{问题一最终测试的模型比较（18 条序列指标取平均）}\label{tab:q1model}
\begin{tabular}{lrrr}\toprule
模型 & MAE/GPU & RMSE/GPU & WAPE\\\midrule
24 小时季节朴素 & 35.3032 & 46.9509 & 1.0809\\
168 小时季节朴素 & 35.9120 & 46.9954 & 1.2171\\
验证集选定模型 & 25.5892 & 34.5070 & 0.8231\\\bottomrule
\end{tabular}
\end{table}

\begin{table}[H]\centering
\caption{验证选模后按任务类型汇总的测试误差}\label{tab:q1type}
\begin{tabular}{lrrr}\toprule
任务类型 & 平均 MAE/GPU & 平均 RMSE/GPU & 平均 WAPE\\\midrule
实时推理 & 3.1140 & 4.3412 & 0.8465\\
批量推理 & 12.7318 & 16.9469 & 0.7328\\
AI 训练 & 60.9217 & 82.2330 & 0.8899\\\bottomrule
\end{tabular}
\end{table}

AI 训练的绝对误差最高，原因是单任务 GPU 规模和尖峰幅度显著大于推理任务；实时推理绝对误差较小，但因真实需求基数较低，WAPE 并未同步降到最低。故三项指标必须联合解释，不能以单一 WAPE 判断容量调度风险。选定的 17 个 Huber 模型覆盖多组 $\epsilon$ 与 $\alpha$：其中 $\epsilon=1.75,\alpha=10^{-4}$ 出现 6 次，其余组合分散出现，说明逐序列验证选择并非退化为单一公共参数。

18 条序列中 17 条选择 Huber，RegionC 的实时推理选择 168 小时季节朴素。相对两条朴素基线，验证集选定模型的平均 MAE 至少下降 27.5\%，说明独立验证选模在本测试窗口内有效；但这只是一次保留测试，不代表长期泛化误差。将 18 条序列聚合后的逐小时预测与真实值见图~\ref{fig:q1forecast}。
\begin{figure}[H]\centering
\includegraphics[width=.86\textwidth]{q1_forecast_actual.png}
\caption{第 2376 至 2399 小时聚合 GPU 需求：预测与真实值}\label{fig:q1forecast}
\end{figure}

\subsubsection{538 个实际任务的末端调度}
最后 24 小时不用预测需求，而使用 538 个实际到达任务。确定性可行构造依次按实时属性、截止、持续时间、GPU 需求和任务编号排序，再按等待、时延和新增峰值利用率选择候选。共有 68 个任务在第 2400 至 2405 小时完成，且第 2406 小时任务占用为零。峰值利用率与按重叠时间折算的平均利用率分别为
\begin{equation}
U^{peak}_{rt}=\frac{\sum_i g_i\chi_{it}\mathbf 1(r_i=r)}{\bar G_r},\qquad
U^{avg}_{rt}=\frac{\sum_i g_i\omega_{it}\mathbf 1(r_i=r)}{\bar G_r}.
\end{equation}
图~\ref{fig:gpu} 来自 solve v12 的逐时图形数据，完整调度表接受任务唯一性、时窗、时延和三类容量全量审计。
\begin{figure}[H]\centering
\includegraphics[width=.86\textwidth]{fig_1_gpu_utilization.png}
\caption{最后 24 小时各区域 GPU 利用率}\label{fig:gpu}
\end{figure}

\subsection{问题二：五方向增量任务搜索}
问题二不启用储能。逐时能源满足
\begin{align}
R^{use}_{rt}&=\min\{R^{av}_{rt},L^F_{rt}\},\\
G^{buy}_{rt}+R^{use}_{rt}&=L^F_{rt},\\
R^{use}_{rt}+R^{sell}_{rt}+R^{curt}_{rt}&=R^{av}_{rt},
\end{align}
且购售电互斥并受进出口上限约束。对某个任务移动，只重算原执行区、新执行区及任务占用小时上的 GPU、IT、设施功率和能源数组；若任一硬边界失败，候选立即淘汰。通过增量检查后，根据搜索方向构造分数：成本方向优先降低净成本，碳方向优先降低购电碳排放，新能源方向优先减少弃电，服务方向优先减少时延与等待，平衡方向同时考虑上述变化。只有严格改善才更新 incumbent。

该增量过程不替代最终验证。每个方向结束后重新从任务表构造全部区域—小时数组，并对 50000 个任务执行完整审计。每个搜索方向以同一可行任务方案为起点，完整评价 10000 个候选；“方向名”仅说明候选排序准则。

\begin{table}[H]\centering
{\small
\caption{问题二五类有限搜索结果}\label{tab:q2schemes}
\begin{tabular}{lrrrrrr}\toprule
方向 & 成本/亿元 & 碳/t & 新能源率 & 时延/ms & 迁移数 & 接受数\\\midrule
成本 & -4.2334 & 10.2851 & 0.680386 & 6.2235 & 368 & 422\\
碳 & -4.1966 & 14.3916 & 0.679078 & 5.1594 & 209 & 0\\
新能源 & -4.2230 & 14.3916 & 0.680193 & 7.0758 & 574 & 761\\
服务 & -4.1966 & 14.3916 & 0.679078 & 5.1594 & 209 & 0\\
平衡 & -4.2180 & 16.5160 & 0.679962 & 6.5824 & 561 & 854\\\bottomrule
\end{tabular}
}
\end{table}

平衡方向被用于后续联合基线。五次搜索共评价 50000 个候选，每个方向均对 50000 个任务执行一次全量审计。平衡方案的服务质量为 0.990393，最大等待为 485 小时；这说明代理评分允许以等待损失换取其他指标，不能只根据“平衡”名称推断服务最优。

\begin{table}[H]\centering
{\small
\caption{问题二成本现金流分解}\label{tab:q2cost}
\begin{tabular}{lrrrrr}\toprule
方向 & 购电/MWh & 售电/MWh & 购电支出/万元 & 售电收入/亿元 & 净成本/亿元\\\midrule
成本 & 25.8055 & 1343742.0294 & 0.8516 & 4.2335 & -4.2334\\
碳 & 40.4664 & 1329591.2956 & 1.2328 & 4.1967 & -4.1966\\
新能源 & 40.4664 & 1339961.0855 & 1.2328 & 4.2231 & -4.2230\\
服务 & 40.4664 & 1329591.2956 & 1.2328 & 4.1967 & -4.1966\\
平衡 & 43.7999 & 1337864.0359 & 1.4424 & 4.2182 & -4.2180\\\bottomrule
\end{tabular}
}
\end{table}

表~\ref{tab:q2cost} 说明负成本几乎完全由售电收入驱动。

平衡方案相对于未搜索基线把净成本由 $-4.196581$ 亿元降至 $-4.218017$ 亿元；碳排放由 14.3916 增至 16.5160 tCO$_2$；它是有限候选集中的折中结果，不构成跨方案 Pareto 最优。

\subsection{问题三：离散 SOC 动态规划与候选筛选}
设施负荷固定为 $L^F_{rt}=\pi_r(L^{BAI}_{rt}+L^N_{rt})$。令 $S_{rt}$ 为时段末 SOC，充放电功率为 $C_{rt},D_{rt}$，则
\begin{align}
S_{rt}&=S_{r,t-1}+\eta_r^cC_{rt}-D_{rt}/\eta_r^d,\\
S_r^{min}\le S_{rt}&\le S_r^{max},\qquad S_{r,2406}=S_r^0.
\end{align}
每区使用 21 个等距状态并额外加入 $S_r^0$。设 $\mathcal S_r$ 为离散状态集合，$\phi_t(s,s')$ 为从 $s$ 转移至 $s'$ 时对应策略的单步代价，则动态规划递推为
\begin{equation}
J_t(s')=\min_{s\in\mathcal S_r:\,(s,s')\ \text{可行}}\{J_{t-1}(s)+\phi_t(s,s')\},
\qquad J_{-1}(S_r^0)=0.
\end{equation}
可行转移由充放电功率、购电上限与能源守恒共同决定；正式方案终点固定为 $S_r^0$。成本策略取购电支出减售电收入，碳策略放大购电碳排放，新能源策略主要惩罚弃电，峰值策略使用购电平方，爬坡策略惩罚相对参考购电的绝对变化，平衡策略组合成本、碳、弃电和购电平方。逐时保留最优前驱并在终点回溯完整 SOC 轨迹。

峰值、爬坡和等效循环定义为
\begin{equation}
P=\sum_r\max\{0,\max_t(G^{buy}_{rt}-G^{sell}_{rt})\},\quad
V=\sum_{r,t}|N_{rt}-N_{r,t-1}|,
\end{equation}
\begin{equation}
N^{cycle}=\sum_r\frac{\sum_t(C_{rt}+D_{rt})}{2(S_r^{max}-S_r^{min})}.
\end{equation}

\begin{table}[H]\centering
{\small
\caption{问题三七策略的经济、碳与新能源指标}\label{tab:q3policies}
\begin{tabular}{lrrrc}\toprule
策略 & 成本/亿元 & 碳/t & 新能源率 & 唯一轨迹\\\midrule
no-action & -4.1982 & 8.8506 & 0.679150 & 是\\
成本 & -4.5875 & 0.0000 & 0.711735 & 是\\
碳 & -4.1990 & 0.0000 & 0.679300 & 是\\
新能源 & -4.5856 & 4.6997 & 0.771475 & 是\\
峰值 & -4.1990 & 0.0000 & 0.679300 & 否（同碳策略）\\
爬坡 & -4.1982 & 11.4975 & 0.679305 & 是\\
平衡 & -4.5858 & 0.2241 & 0.771439 & 是\\\bottomrule
\end{tabular}
}
\end{table}

\begin{table}[H]\centering
{\small
\caption{问题三七策略的峰值、爬坡与循环指标}\label{tab:q3operation}
\begin{tabular}{lrrrr}\toprule
策略 & 峰值/MW & 爬坡/MW & 等效循环 & 候选内非支配\\\midrule
no-action & 31.5979 & 63.1957 & 0.000 & 否\\
成本 & 0.0000 & 0.0000 & 240.186 & 是\\
碳 & 0.0000 & 0.0000 & 2.146 & 否\\
新能源 & 2.7455 & 15.2020 & 1568.882 & 是\\
峰值 & 0.0000 & 0.0000 & 2.146 & 否\\
爬坡 & 40.9020 & 81.8040 & 2.496 & 否\\
平衡 & 0.3740 & 0.7480 & 1568.242 & 是\\\bottomrule
\end{tabular}
}
\end{table}

选中的平衡策略相对 no-action：成本减少 38764304.59 元，碳排放减少 8.6265 tCO$_2$，峰值减少 31.2239 MW，绝对爬坡减少 62.4478 MW。累计充、放电分别为 906765.54 MWh 和 786695.83 MWh。图~\ref{fig:soc} 显示终端 SOC 回到初始值。
\begin{figure}[H]\centering
\includegraphics[width=.86\textwidth]{fig_2_storage_soc.png}
\caption{问题三平衡策略的六区域 SOC 轨迹}\label{fig:soc}
\end{figure}

循环次数很高，主要源于未计退化成本和有限 SOC 网格允许逐日重复利用价差与新能源富余。它不违反题面给定硬约束，但会削弱工程可实施性；因此本文只把该轨迹称为候选集内非支配方案。终端 SOC 放宽实验下，各区终端值降至 35、32、30、90、82、85 MWh，而成本最小策略的成本与强制回到初值时相同，说明当前成本最小候选的结果并非由终端库存被迫抬高造成。

\subsection{问题四：任务—储能动态交替搜索}
联合模型用真实任务轨迹重建 $L^F_{rt}$，并继承前三问全部任务、能源和 SOC 约束。联合基线从问题二平衡任务方案出发，储能侧仍从真实初始 SOC 重启。候选集内的基准联合结果为：成本 $-458716347.97$ 元，碳排放 0.268398 tCO$_2$，网络时延 6.582432 ms，服务质量 0.990393，新能源利用率 0.771883，峰值购电 0.364428 MW。

每轮联合评价使用七维向量
\begin{equation}
\mathbf z=(C,E,-U,P,V,L_{net},\overline W),
\qquad
\Psi(\mathbf z)=\frac{1}{7}\sum_{j=1}^{7}\frac{z_j-z_j^0}{\max(|z_j^0|,1)}.
\end{equation}
若设置碳目标且候选超过目标，程序按超额比例加入显式罚项。这样可避免亿元级成本在未经尺度化时压倒碳、峰值和服务指标；同时，罚项只用于声明启发式的候选排序，不能转化为不可行性证明。

实际交替过程如下。
\begin{enumerate}
\item 对当前任务方案重建逐时设施负荷，并求离散 SOC 储能策略；
\item 由当前电价、碳强度、购电和弃电形成任务移动信号；
\item 按固定次序评价 600 个任务移动候选，只接受通过增量边界的严格改进；
\item 对候选执行 50000 任务全量审计，若接受则以新任务与新储能作为下一轮 incumbent；
\item 达到 8 轮上限或连续 2 轮无改善时停止；每个情景独立重启。
\end{enumerate}

基准情景完成 2 轮、评价 1200 个候选且未接受轮次；10\% 碳削减情景在第 2 轮接受一次改进，总计 4 轮和 2400 个候选；种子 2028 的波动情景在第 1 轮接受一次改进，总计 3 轮和 1800 个候选。所有情景合计 37 轮、22200 个候选。


轮次审计进一步区分“局部移动改善”和“整轮联合接受”：基准两轮分别出现 44 和 53 次局部改善，但储能重建后的联合分数均变差，故不更新 incumbent；碳削减 10\% 情景仅第 2 轮被接受，联合分数由 58.075356 降至 0.101290，随后连续两轮无改善停止；波动种子 2028 仅第 1 轮被接受，分数由 303.034335 降至 294.820472。该对照表明邻域随前一轮储能信号动态重建，而不是从固定候选表中挑选。完整逐轮轨迹保留在结构化证据中。

\begin{table}[H]\centering
{\small
\caption{基准与碳约束情景}\label{tab:q4scenarios}
\begin{tabular}{lrrrrrr}\toprule
水平 & 轮数 & 迁移数 & 成本/亿元 & 碳/t & 新能源率 & 峰值/MW\\\midrule
基准 & 2 & 561 & -4.583419 & 0.5789 & 0.768912 & 0.4438\\
削减 0.00 & 2 & 561 & -4.583419 & 0.5789 & 0.768912 & 0.4438\\
削减 0.10 & 4 & 582 & -4.583651 & 0.4245 & 0.768957 & 0.4438\\
削减 0.20 & 2 & 561 & -4.583418 & 0.5789 & 0.768912 & 0.4438\\
削减 0.30 & 2 & 561 & -4.583417 & 0.5789 & 0.768912 & 0.4438\\\bottomrule
\end{tabular}
}
\end{table}

\begin{table}[H]\centering
{\small
\caption{电价与售电机制情景}\label{tab:q4price}
\begin{tabular}{llrrrrr}\toprule
类型 & 系数 & 轮数 & 迁移数 & 成本/亿元 & 新能源率 & 峰值/MW\\\midrule
峰谷价差 & 0.75 & 2 & 561 & -4.583419 & 0.768912 & 0.4438\\
峰谷价差 & 1.00 & 2 & 561 & -4.583419 & 0.768912 & 0.4438\\
峰谷价差 & 1.25 & 2 & 561 & -4.583419 & 0.768912 & 0.4438\\
售电机制 & 0.00 & 2 & 561 & 0.000005 & 0.768926 & 0.4438\\
售电机制 & 0.50 & 2 & 561 & -2.291688 & 0.768917 & 0.4438\\
售电机制 & 1.00 & 2 & 561 & -4.583419 & 0.768912 & 0.4438\\\bottomrule
\end{tabular}
}
\end{table}

\begin{table}[H]\centering
{\small
\caption{新能源水平与波动情景}\label{tab:q4renewable}
\begin{tabular}{llrrrrr}\toprule
类型 & 水平/种子 & 迁移数 & 成本/亿元 & 碳/t & 新能源率 & 峰值/MW\\\midrule
水平 & 0.80 & 561 & -3.914911 & 500.5503 & 0.892089 & 92.1180\\
水平 & 1.00 & 561 & -4.583419 & 0.5789 & 0.768912 & 0.4438\\
水平 & 1.20 & 561 & -4.592361 & 0.0000 & 0.652385 & 0.0000\\
波动 & 2026 & 561 & -4.534058 & 760.7500 & 0.768841 & 273.9082\\
波动 & 2027 & 561 & -4.532170 & 852.0600 & 0.768729 & 298.6499\\
波动 & 2028 & 583 & -4.537630 & 464.9554 & 0.768887 & 146.1294\\\bottomrule
\end{tabular}
}
\end{table}

20\% 与 30\% 碳削减目标的状态是 \texttt{target\_infeasible\_within\_declared\_heuristic}，即声明搜索内未达到。售电系数改变几乎只影响净成本，新能源水平和波动则同时改变碳排放与峰值购电，说明新能源不确定性是联合运行的主要敏感源。

\subsection{统一约束审计与方法边界}
任务审计逐项检查唯一执行、实时即时开工、到达、截止、单向时延、第 2406 小时禁占用、GPU、IT 和设施功率。能源审计检查
\begin{align}
e^{energy}_{rt}&=G^{buy}_{rt}+R^{av}_{rt}+D_{rt}-L^F_{rt}-C_{rt}-G^{sell}_{rt}-R^{curt}_{rt},\\
e^{SOC}_{rt}&=S_{rt}-S_{r,t-1}-\eta_r^cC_{rt}+D_{rt}/\eta_r^d.
\end{align}
只有最大残差、互斥、功率边界与终端等式全部通过才写出正式结果。

% MMW-LIMITATION: D1
\paragraph{实现边界 D1}Q2 与 Q4 是预声明候选预算内的 best-found 可行解，不是穷举最优解。
% MMW-LIMITATION: D2
\paragraph{实现边界 D2}Q3 使用 21 个 SOC 网格状态，未计题面未给出的退化成本与循环寿命上限。
% MMW-LIMITATION: D3
\paragraph{实现边界 D3}Q4 的随机波动只使用固定种子 2026、2027、2028，并逐情景独立重启。
% MMW-LIMITATION: L1
\paragraph{证据边界 L1}无独立 Oracle，最高可信等级为 scenario-feasible。
% MMW-LIMITATION: L2
\paragraph{证据边界 L2}候选集内非支配不等于完整决策空间的 Pareto 前沿。
% MMW-LIMITATION: L3
\paragraph{证据边界 L3}Q2 方案名表示搜索方向，不表示跨方案的对应目标最小值。
